\documentclass[11pt]{article}
\usepackage[margin=1in]{geometry}
\usepackage{amsmath,amssymb}
\usepackage{array}
\usepackage{booktabs}
\usepackage{hyperref}

\title{M.I.N.D. v3 Research Plan: Dynamic Graph-Calibrated EEG Decoding for Reliable Motor-Imagery BCIs}
\author{JR-Academy 6857}
\date{}

\begin{document}
\maketitle

\begin{abstract}
M.I.N.D. v1 established a reliability-focused classical motor-imagery EEG
pipeline, and M.I.N.D. v2 added modern deep backbones while preserving strict
cross-session evaluation and selective prediction. M.I.N.D. v3 should move
beyond using known architectures as comparisons and propose a new reliability-
aware architecture. This document outlines a research plan for a dynamic
graph-calibrated EEG decoder that jointly models channel relationships,
temporal motor-imagery structure, and uncertainty under cross-session shift.
The plan is intentionally provisional: the architecture must be checked against
the literature, implemented carefully, and evaluated against v1 and v2 before
any novelty or performance claims are made.
\end{abstract}

\section{Motivation}

The central research claim of M.I.N.D. is that motor-imagery BCIs should be
evaluated as reliability-aware controllers, not ordinary forced-choice
classifiers. In realistic assistive use, the system should act when neural
evidence is reliable and abstain when uncertainty is high.

M.I.N.D. v2 shows that modern deep models improve performance over the v1
classical LDA baseline. However, v2 is not architecturally novel: EEGNet is an
established model, and the spatial-attention backbone is a lightweight
extension rather than a new decoding paradigm. The goal of v3 is to make the
model itself reliability-aware.

\section{Proposed v3 Architecture}

The proposed v3 model is a \textbf{Dynamic Graph-Calibrated EEG Transformer}
for motor-imagery EEG. The core idea is to represent EEG channels as graph
nodes, learn trial-specific or session-specific channel relationships, process
the resulting graph features over time, and train the probability output to be
selectively reliable.

\subsection{Input}

Each EEG trial is represented as
\[
  X \in \mathbb{R}^{C \times T},
\]
where \(C\) is the number of EEG channels and \(T\) is the number of time
samples. Euclidean alignment remains an optional label-free preprocessing
step, preserving the v1/v2 cross-session adaptation philosophy.

\subsection{Dynamic Channel Graph}

Each channel is treated as a node. Instead of using a fixed adjacency matrix,
v3 should explore a dynamic adjacency matrix
\[
  A(X) \in \mathbb{R}^{C \times C},
\]
estimated from the trial, the training session, or a learned attention
mechanism. Candidate graph definitions include:
\begin{itemize}
  \item covariance similarity between channels,
  \item phase-locking or coherence-based connectivity,
  \item learned query-key attention over channel embeddings,
  \item a hybrid graph combining electrode geometry with learned functional
  connectivity.
\end{itemize}
The literature review must determine which graph construction is justified and
whether similar dynamic EEG graph models already exist.

\subsection{Graph Spatial Encoder}

A lightweight graph attention or graph convolution block maps raw channel
features into spatially contextualized embeddings:
\[
  H = \mathrm{GraphEncoder}(X, A(X)).
\]
The purpose is to model relationships among electrodes directly, rather than
forcing the network to infer all channel interactions through ordinary spatial
convolutions.

\subsection{Temporal Encoder}

After spatial graph encoding, the model applies a temporal module:
\[
  Z = \mathrm{TemporalEncoder}(H).
\]
Candidate temporal encoders include:
\begin{itemize}
  \item depthwise temporal convolutions,
  \item temporal convolutional networks,
  \item lightweight transformer encoders,
  \item conformer-style local convolution plus attention.
\end{itemize}
The initial implementation should favor the simplest version that improves
over v2 without becoming unstable on small subject-specific training sets.

\subsection{Reliability-Aware Output Head}

The key novelty should not stop at the graph backbone. v3 should include an
output head and training objective designed for calibrated selective
prediction. Candidate heads include:
\begin{itemize}
  \item temperature-scaled softmax trained with cross-entropy plus Brier loss,
  \item evidential or Dirichlet classification head,
  \item auxiliary uncertainty head predicting whether a trial should be
  accepted at the target coverage,
  \item selective-classification loss that jointly optimizes classification
  risk and coverage.
\end{itemize}
The safest first version is cross-entropy plus Brier loss, followed by
post-hoc temperature scaling. More experimental uncertainty heads should be
treated as ablations until validated.

\section{Research Questions}

\begin{enumerate}
  \item Does a dynamic EEG channel graph improve cross-session motor-imagery
  decoding over EEGNet and the v2 spatial-attention model?
  \item Does reliability-aware training reduce ECE and Brier score compared
  with post-hoc calibration alone?
  \item Does v3 improve selective accuracy and risk-coverage behavior at
  matched coverage?
  \item Are gains consistent across subjects and across external MOABB
  datasets?
  \item Which part matters most: Euclidean alignment, dynamic graph modeling,
  temporal attention, calibration-aware loss, or selective prediction?
\end{enumerate}

\section{Required Literature Review}

Before implementation, the following areas must be reviewed carefully:
\begin{itemize}
  \item EEGNet and compact convolutional EEG models,
  \item graph neural networks for EEG classification,
  \item dynamic graph construction for EEG or biosignals,
  \item EEG transformers and conformer-style models,
  \item calibration methods for neural networks,
  \item selective classification and risk-coverage evaluation,
  \item evidential deep learning and uncertainty estimation,
  \item Euclidean alignment and Riemannian/domain-adaptation methods for EEG.
\end{itemize}
The final paper must explicitly distinguish what is borrowed from prior work
and what is new in M.I.N.D. v3.

\section{Evaluation Protocol}

v3 should preserve the strict leakage controls from v1 and v2:
\begin{itemize}
  \item train on \(A0xT\) and evaluate once on \(A0xE\),
  \item fit Euclidean alignment without labels,
  \item use only training-session validation data for calibration, threshold
  selection, early stopping, and ensemble/model selection,
  \item never use evaluation labels until final metric computation.
\end{itemize}

The required comparisons are:
\begin{itemize}
  \item v1 calibrated LDA,
  \item v2 EEGNet calibrated,
  \item v2 spatial-attention calibrated,
  \item v2 validation-weighted deep ensemble,
  \item v3 graph-calibrated model,
  \item v3 ablations.
\end{itemize}

\section{Ablation Plan}

The v3 ablation table should include:
\begin{itemize}
  \item static graph vs dynamic graph,
  \item geometry-only graph vs learned functional graph,
  \item graph encoder only vs graph plus temporal encoder,
  \item cross-entropy only vs cross-entropy plus Brier loss,
  \item with vs without temperature scaling,
  \item max-prob vs entropy vs margin uncertainty score,
  \item with vs without Euclidean alignment.
\end{itemize}

\section{Metrics}

The v3 evaluation should report the same metrics as v2:
\begin{itemize}
  \item all-trial accuracy,
  \item macro F1,
  \item multiclass Brier score,
  \item Expected Calibration Error,
  \item selective accuracy at 40\%, 60\%, and 80\% coverage,
  \item selective risk at matched coverage,
  \item full risk-coverage curves,
  \item trial-weighted reliability diagrams with bin support.
\end{itemize}

\section{Expected Contribution}

If validated, M.I.N.D. v3 could make a stronger novelty claim than v2:
\begin{quote}
M.I.N.D. v3 introduces a dynamic graph-calibrated EEG decoder that learns
trial-specific channel relationships and optimizes probability reliability for
selective motor-imagery BCI control under cross-session shift.
\end{quote}

This claim should only be made if the literature review confirms that the
combination of dynamic graph modeling, calibration-aware learning, and
selective-prediction evaluation is meaningfully distinct from prior EEG
decoding work.

\section{Risks and Failure Modes}

The main risks are:
\begin{itemize}
  \item the dynamic graph model may overfit because BCI IV 2a is small,
  \item reliability-aware losses may improve calibration but reduce accuracy,
  \item graph construction may add complexity without improving subject-level
  consistency,
  \item similar architectures may already exist in the EEG literature,
  \item external MOABB validation may weaken the claim if gains do not
  transfer.
\end{itemize}

The fallback plan is to keep v3 modular: if the full graph transformer is too
unstable, evaluate a smaller graph-attention spatial encoder with the same
v2 reliability framework.

\section{Conclusion}

M.I.N.D. v3 should be treated as a research project, not just an implementation
upgrade. The strongest path is to preserve the leakage-aware reliability
evaluation from v1 and v2 while proposing a model whose structure directly
matches the scientific problem: EEG channels interact dynamically, session
shift changes those relationships, and assistive BCIs need calibrated
confidence rather than forced predictions on every trial.

\end{document}
